\section{Research Goal}
\label{sec:research_objectives_research_goal}

The main research goals are:

\begin{enumerate}
    \item \textbf{Unified Framework for Semantic Compositionality:} 
    To develop a unified theoretical and empirical framework that explains 
    when and why pretrained generative models---whether energy-based product of experts, mixtures or density-based (e.g., SuperDiff)---yield 
    \emph{semantically meaningful and plausible} compositions across 
    different domains.  
    This includes systematically characterising the behaviour, success 
    modes, and failure modes of existing compositional operators and 
    identifying the regimes in which they respect the underlying 
    generative factors.

    \item \textbf{Robust Property-Guided Compositional Diffusion:}
    To extend compositional generative modelling to settings in which 
    continuous properties (e.g., severity, spatial attributes, style 
    strength) guide the sampling process on top of composite densities.  
    This goal further includes developing or adapting 
    context-aware guidance methods (such as context-guided diffusion) 
    that maintain calibrated, uncertainty-aware behaviour under 
    distribution shift, thereby improving the robustness and semantic 
    reliability of guided compositions.

    \item \textbf{Intrinsic Geometric Preconditions for Compositional Correctness:}
    To identify and formalise the \emph{geometric} conditions under which
    pretrained diffusion models can be composed coherently. This goal focuses on
    characterising how curvature, metric structure and manifold compatibility influence the feasibility of energy-based product of experts, mixtures
    and density-based compositions. We further aim to determine when the
    manifolds underlying different diffusion models intersect smoothly or conflict geometrically, and how
    such properties predict the success or failure of compositional operators
    across domains.
\end{enumerate}